
\section{Young Graph by Alcove Geometry}
This section is a small note, devoted to the idea of realising the Young Graph by alcove geometry. Firstly, consider the type $B_{\infty}$ Weyl group, $G$:
\[ G=\langle s_0, s_1,\dots \ | \ s_k^2=(s_0 s_1)^4=(s_{l}s_{l+1})^3=(s_i s_j)^2=e \ \text{for }l>0, |i-j|>1 \rangle,\]
acting on $\R^\infty$ (the space of finite real sequences) as follows:
\begin{align*}
    s_0(x_1,x_2,\dots,x_n)&=(-x_1,x_2,\dots,x_n),\\
    s_i(x_1,\dots,x_i,x_{i+1},\dots,x_n)&=(x_1,\dots,x_{i+1},x_i,\dots,x_n),
\end{align*}
where we use $(x_1,x_2,\dots,x_n)$ as short hand for the same infinite sequence, padded by 0s on the right. By allowing the $x_k$ to possibly be 0, we can ensure the last definition is always sensible. We can consider the $s_0$ as acting by reflection in the hyperplane defined by $x_1=0$, and the $s_i$ as acting by reflection in the hyperplane defined by $x_i=x_{i+1}$.

Given an element $(x)\in \R^\infty$ we can apply $s_0$ (conjugated by permutations to act in each component) so that all non-zero elements are positive, and then apply a permutation to order all entries in descending order. Thus a fundamental domain, $\mathcal{D}$, for $G\circlearrowright \R^\infty$, is the set of finite, descending, positive sequences:
\[\mathcal{D}= \left\{(x_1,x_2,\dots, x_n,0,\dots) \in \R^\infty \ \Bigg| \ \begin{matrix}\text{ for some } n\in \N \text{ such that }\\ x_1\geq x_2\geq \dots \geq x_n>0\end{matrix} \right\}.\]
Now, let us augment our group to be $\hat{G}$ where we add another generator $r_0$ to $G$ which satisfies the exact same relations as $r_0$ and has no ``braiding" relation with $s_0$. We give a Dynkin diagram for the group $\hat{G}$ below:
\[ \begin{tikzcd}[column sep=tiny,row sep = tiny]
    s_0\ar[dr,equal]\\
    &s_1\ar[r,dash]&s_2\ar[r,dash]&s_3\ar[r,dash]&\dots\\
    r_0\ar[ur,equal]\ar[uu,dash,"\infty"]
\end{tikzcd}\]
Now let $r_0$ act on $\R^\infty$ by the (affine) reflection in the hyperplane defined by $x_1=1$, that is,
\[r_0(x_1,x_2,\dots,x_n)=(2-x_1,x_2,\dots,x_n).\]
We then observe that $r_0 s_0$ acts by $x_1\mapsto x_1+2$ (fixing all other components). For a partition, $\l\vdash n$, we may view $\l=(\l_1,\l_2,\dots, \l_r)\in \mathcal{D}$ with $\l_i\geq \l_{i+1}$ for all $i$. Now define
\[\mathcal{D}_{\l}:=\{x\in \R^\infty \ | \ x-\l \in [0,1)^\infty  \}.\]
Clearly $\mathcal{D}=\sqcup_{\l} \mathcal{D}_{\l}$ and furthermore, we claim that each $\mathcal{D}_{\l}$ is a fundamental domain for $\hat{G}\circlearrowright \R^\infty$ (or an alcove) which shares a face with $\mathcal{D}_{\l'}$, if and only if, $\l'$ is a partition obtained from $\l$ by ``adding or removing a box". Suppose that $\l$ and $\l'=(\l_1,\dots,\l_i+1,\dots ,\l_r)$ are both partitions. Then $\mathcal{D}_{\l}$ and $\mathcal{D}_{\l'}$ (or rather, their closures) share the face
\[\overline{\mathcal{D}_{\l}}\cap\overline{\mathcal{D}_{\l'}}=\left\{\begin{matrix}(\l_1+\xi_1,\dots,\l_i+1,\dots,\l_r+\xi_r,\xi_{r+1},\dots,\xi_k)\in \mathcal{D}  \\ :\text{for some } k \in \N, \text{ where } \xi_l\in [0,1]\end{matrix} \right\}.\]
Therefore, by identifying each alcove $\mathcal{D}_{\l}$ with the underlying partition $\l$, we obtain the Young-graph as the dual graph of this alcove geometry.


